%! Sinusoidal Function Context and Data Modeling — Unit 6, Lesson 6.8 — Exit Ticket
\documentclass[10pt]{article}
\usepackage{precalculus-article}
\usepackage{precalculus-boxes}
\newcommand{\TallMath}[1]{$\displaystyle #1\rule[-1.4em]{0pt}{3.2em}$}

\begin{document}

\pageheader{Unit 6, Lesson 6.8}{Exit Ticket}
\namedateperiod

\vspace{0.12in}

\begin{notesbox}{Exit Ticket --- Answer independently}

The table below gives values of a periodic function $f$ over one period.

\vspace{6pt}
\begin{center}
\renewcommand{\arraystretch}{1.4}
\begin{tabular}{|c|c|c|c|c|c|c|}
\hline
$x$ & 0 & 2 & 4 & 6 & 8 & 10 \\
\hline
$f(x)$ & 3 & 11 & 3 & $-5$ & 3 & 11 \\
\hline
\end{tabular}
\end{center}

\vspace{8pt}
\textbf{1.}\quad Identify the maximum and minimum values of $f$.

Maximum $= $ \underline{\hspace{2cm}} \qquad Minimum $= $ \underline{\hspace{2cm}}

\vspace{8pt}
\textbf{2.}\quad Compute the amplitude and midline.

$A = \dfrac{\max - \min}{2} = $ \underline{\hspace{3.5cm}}

\vspace{4pt}
$D = \dfrac{\max + \min}{2} = $ \underline{\hspace{3.5cm}}

\vspace{8pt}
\textbf{3.}\quad Estimate the period $k$ from the table. Explain how you found it.

$k = $ \underline{\hspace{2cm}} \quad Explanation: \writeline

\vspace{8pt}
\textbf{4.}\quad Using your period, compute $B = \dfrac{2\pi}{k} = $ \underline{\hspace{3.5cm}}

\vspace{8pt}
\textbf{5.}\quad A student uses the model $f(x) = 8\cos\!\left(\dfrac{\pi}{4}(x - 2)\right) + 3$
to predict $f(3)$. Is $x = 3$ inside the contextual domain for this data set?
Circle one: \quad \textbf{YES} \quad / \quad \textbf{NO}

\vspace{4pt}
Explain: \writelines{2}

\end{notesbox}

\end{document}
